% §16.2 Varying Rates Over Time
%
% Monograph register: phenomenon (fpr(t)), which axiom (Axiom 2
% stability), consequences for iid tools, Azuma-Hoeffding, Bloom-past-
% capacity worked example, time-sorted diagnostic. Folds salvaged ex-5
% as rem:composition-jensen (random stage rates make the composition
% theorem optimistic by Jensen). Cross-refs ax:element-indep,
% ax:block-rates-indep, sec:axiom-2, thm:chernoff, sec:bloom-from-scratch,
% thm:composition, sec:where-framework-helps, sec:correlated-noise,
% sec:adversarial-inputs. DEFINED: sec:varying-rates, prop:azuma-hoeffding,
% rem:composition-jensen.

\section{Varying Rates Over Time}
\label{sec:varying-rates}

The first failure mode is the most common in practice. The error rate is
not a fixed property of the channel; it changes with the query index.
Where the earlier chapters wrote $\fprate$ as a scalar, this section
writes $\fprate(t)$ and asks what happens to the iid analysis when the
scalar reading is no longer available.

\subsection*{The phenomenon}

Three examples recur through the chapter. A \emph{drifting classifier},
trained on one distribution and deployed on another, sees its per-query
false-positive rate climb as the deployment distribution moves away from
the training one, from $\fprate(t) \approx 0.02$ early to $0.15$ months
later. A \emph{Bloom filter past capacity}, sized for $n = 1000$ at
$\fprate \approx 0.01$ but loaded with $5000$ items, watches its fill
ratio rise from $30\%$ toward $100\%$ and its false-positive rate climb
from $0.01$ to about $0.5$, following $\fprate(t) \approx (1 -
e^{-kt/m})^k$ from \Cref{sec:bloom-from-scratch} with $t$ the insertion
count at query time. A \emph{drifting sensor} sees its error rate track
its operating temperature over the day. What unites the three is
structural: the rate is a non-constant function of the query index, and
treating it as constant produces bounds that underestimate the actual
error frequency, sometimes by large factors.

\subsection*{Which axiom is violated}

\Cref{sec:axiom-2} stated Axiom 2 in two equivalent forms, a conditional
independence of the per-block rates and a stability of the channel matrix
$Q$ across queries. The stability form breaks here: the per-query rate at
step $t$ is $\fprate(t)$, not $\fprate$, so the joint distribution over a
query sequence no longer factors through a single $Q$. The practical
signature is that the same query, repeated at different times, draws
different rate parameters. Time-varying rates are compatible with
\Cref{ax:element-indep} (Axiom 1) in many cases: if the drift is governed
by an exogenous time variable and errors at distinct queries remain
conditionally independent given the drift state, only Axiom 2 fails. The
two axioms are logically independent.

\subsection*{Consequences for the iid tools}

Three workhorses break. The \emph{composition theorem of chapter 3}
($\eta_{\text{total}} = 1 - \prod_i (1 - \eta_i)$) assumed each $\eta_i$
fixed; when the $\eta_i$ are time-indexed or random, the product is an
average over a trajectory, and the error along the worst trajectory can
exceed it. The \emph{PPV formula of chapter 7} ($\PPV = \pi \tprate /
(\pi \tprate + (1 - \pi) \fprate)$) still holds at any single query (the
\Cref{sec:where-framework-helps} salvage) but the time-average PPV need
not equal the PPV at the time-average rates, because PPV is nonlinear in
its arguments. The \emph{Chernoff bound of chapter 11}
(\Cref{thm:chernoff}) assumed iid Bernoulli summands; with varying
$\fprate(t)$ the $X_t$ are independent but not identically distributed,
and the multiplicative form that relied on a single $p$ no longer applies
directly. Its martingale generalization does.

\begin{remark}[Random rates make composition optimistic]
\label{rem:composition-jensen}
The composition formula does not merely lose exactness under random
stage rates; it errs in a fixed direction. If the $\eta_i$ are
independent with means $\mu_i = \mathbb{E}[\eta_i]$, then since $\mathbb{E}[1 - \eta_i] =
1 - \mu_i$ and $-\log(1 - x)$ is convex, Jensen's inequality gives $\mathbb{E}[\,1
- \prod_i (1 - \eta_i)\,] \ge 1 - \prod_i (1 - \mu_i)$, with equality only
when every variance vanishes. Plugging the mean rates into the
deterministic formula therefore \emph{under}-estimates the expected
combined error. For five stages at $\mu = 0.1$ and $\sigma = 0.02$ the
gap is a fraction of a percentage point, but it grows with the stage
count and the rate variance, and it is always in the optimistic
direction.
\end{remark}

\subsection*{The right tool: martingale concentration}

The Azuma-Hoeffding inequality replaces the Chernoff bound off the iid
path. Where Chernoff bounded sums of iid Bernoulli variables,
Azuma-Hoeffding bounds martingale differences that are bounded in size
but otherwise unrestricted in distribution.

\begin{proposition}[Azuma-Hoeffding]
\label{prop:azuma-hoeffding}
Let $\{Y_t\}_{t = 0, 1, \ldots, n}$ be a martingale with $Y_0 = \mathbb{E}[Y_n]$
and bounded differences $|Y_t - Y_{t - 1}| \le c_t$ for each $t \ge 1$.
Then for any $a > 0$,
\[
  \Pr[|Y_n - Y_0| \ge a]
  \;\le\;
  2 \exp\!\left( -\frac{a^2}{2 \sum_{t = 1}^n c_t^2} \right).
\]
\end{proposition}

The inequality is Hoeffding's (1963), with the martingale form from Azuma
(1967); the MGF-plus-Markov argument of chapter 11 carries over once the
per-step MGFs are bounded, and the proof never needs identical
distributions, only the per-step jump bound. Applied to time-varying
false-positive rates, let $X_t \in \{0, 1\}$ indicate a false positive at
query $t$ with $\Pr[X_t = 1 \mid \mathcal{F}_{t-1}] = \fprate(t)$, and set
$Y_t = \sum_{s \le t} (X_s - \fprate(s))$, a martingale with $|Y_t -
Y_{t-1}| \le 1$. With $c_t = 1$,
\[
  \Pr\!\left[\,\Bigl|\,\textstyle\sum_{t=1}^n X_t - \sum_{t=1}^n \fprate(t)\,\Bigr| \ge a\,\right]
  \;\le\;
  2 \exp\!\left( -\frac{a^2}{2 n} \right),
\]
the non-iid generalization of the Chernoff bound: the total false-positive
count concentrates around $\sum_t \fprate(t)$ with Gaussian-shaped tails
of width $O(\sqrt{n})$ even when $\fprate(t)$ is not constant.

\subsection*{Worked example: Bloom filter past capacity}

Take the over-filled filter above, queried $n = 1000$ times after the
fill with non-member inputs, the per-query rate having grown from $0.01$
to $0.5$ along the way. Treating the rate as a constant $\bar\fprate
\approx 0.25$ and applying the iid Chernoff bound to the expected count
$250$ gives $\Pr[X \ge 275] \le e^{-0.83} \approx 0.44$, a useless bound,
because the constant-rate premise is wrong and the mean is misspecified.
The Azuma-Hoeffding bound assumes no constant rate: with $c_t = 1$,
$\Pr[|\sum_t X_t - \sum_t \fprate(t)| \ge a] \le 2 e^{-a^2/2000}$, and
$a = 100$ gives $2 e^{-5} \approx 0.013$. The count lies within $\pm 100$
of the honest trajectory expectation $\sum_t \fprate(t) \approx 230$ with
probability at least $0.987$. The iid analysis fails not because its
bound is weak but because its model (constant rate) does not match the
channel; the martingale analysis succeeds by acknowledging the
non-stationarity.

\subsection*{Diagnostic}

To decide between the two, sort the query log by time and test whether
$\hat\fprate$ is stable across time-ordered partitions: split into early,
middle, and late buckets and compare the three estimates. If they lie
within their sampling intervals, the iid assumption is consistent and
Chernoff applies; if they differ by more, stationarity is broken and
Azuma-Hoeffding along the realized trajectory is the right tool. The test
is cheap and catches the Axiom 2 failure specifically; it says nothing
about the correlated-noise mode of \Cref{sec:correlated-noise} or the
adversarial mode of \Cref{sec:adversarial-inputs}, each of which needs
its own check.
